\documentclass[12pt]{article}
\usepackage[T1]{fontenc}
\usepackage[utf8]{inputenc}
\usepackage{lmodern}
\usepackage{textcomp}
\usepackage{enumitem}
\usepackage{geometry}
\geometry{margin=1in}
\begin{document}
\begin{center}
Answer Key - Health Sciences - K-12 Level Assessment 3 \
\small (Answers are ordered exactly as in the question paper.)
\end{center}

\section*{Multiple Choice}
\begin{enumerate}[label=\arabic*.]
\item \textbf{Answer:} C
\\ \textit{Options:} A) 0.900 g/cm 3; B) 0.730 g/cm 3; C) 1.100 g/cm 3; D) None of the above
\\ \textit{Note:} Correct option: C - 1.100 g/cm 3
\item \textbf{Answer:} D
\\ \textit{Options:} A) 18 carbon atoms with one carbon-carbon double bond in the cis configuration; B) 20 carbon atoms with at least two carbon-carbon double bonds in the cis configuration; C) 18 carbon atoms with at least two carbon-carbon double bonds in the trans configuration; D) 18 carbon atoms with at least two carbon-carbon double bonds in the cis configuration
\\ \textit{Note:} Correct option: D - 18 carbon atoms with at least two carbon-carbon double bonds in the cis configuration
\item \textbf{Answer:} C
\\ \textit{Options:} A) is substantially higher for carbohydrate than for protein; B) is accompanied by a slight decrease in body core temperature.; C) is partly related to sympathetic activity stimulation in the postprandial phase; D) is not attenuated by food malabsorption.
\\ \textit{Note:} Correct option: C - is partly related to sympathetic activity stimulation in the postprandial phase
\item \textbf{Answer:} C
\\ \textit{Options:} A) A reduction in lean body mass; B) A reduction in bone density; C) An increased appetite; D) Impaired immune function
\\ \textit{Note:} Correct option: C - An increased appetite
\item \textbf{Answer:} C
\\ \textit{Options:} A) Alcohol dehydrogenase (ADH); B) Catalase; C) Alcohol dehydrogenase (ADH), cytochrome P 450 2 E 1 (CYP 2 EI), catalase; D) cytochrome P 450 2 E 1 (CYP 2 EI)
\\ \textit{Note:} Correct option: C - Alcohol dehydrogenase (ADH), cytochrome P 450 2 E 1 (CYP 2 EI), catalase
\end{enumerate}

\section*{True / False}
\begin{enumerate}[label=\arabic*.]
\setcounter{enumi}{5}
\item \textbf{Answer:} True
\\ \textit{Options:} A) True; B) False
\\ \textit{Note:} LLM-generated true statement based on MCQ.
\item \textbf{Answer:} True
\\ \textit{Options:} A) True; B) False
\\ \textit{Note:} LLM-generated true statement based on MCQ.
\item \textbf{Answer:} False
\\ \textit{Options:} A) True; B) False
\\ \textit{Note:} LLM-generated false statement. Correct answer: 'Stimulate production of insulin-like growth hormone'.
\item \textbf{Answer:} True
\\ \textit{Options:} A) True; B) False
\\ \textit{Note:} LLM-generated true statement based on MCQ.
\item \textbf{Answer:} False
\\ \textit{Options:} A) True; B) False
\\ \textit{Note:} LLM-generated false statement. Correct answer: 'Result in bone mineral loss from the skeleton'.
\end{enumerate}

\section*{Long Form Response}
\begin{enumerate}[label=\arabic*.]
\setcounter{enumi}{10}
\item \textbf{Answer:} Pepsinogen and gastric lipase are two enzymes secreted into the lumen of the stomach. Explain why this is the correct answer and provide supporting evidence. Reference key facts or concepts that support this choice.
\\ \textit{Note:} Long-form derived from MCQ; award credit for stating the correct idea and giving supporting reasoning.
\item \textbf{Answer:} Enteropeptidase converts trypsinogen to trypsin by cleavage of a peptide sequence that blocks the active site of trypsin.. Explain why this is the correct answer and provide supporting evidence. Reference key facts or concepts that support this choice.
\\ \textit{Note:} Long-form derived from MCQ; award credit for stating the correct idea and giving supporting reasoning.
\end{enumerate}

\vfill
\noindent\textit{End of Answer Key}
\end{document}